\scanpage{44}
\uline{Remarque} On voit facilement par des exemples dans le Th. et les deux corollaires, que les conditions de \unsure{platitudes|platitude} et de \uline{propreté} sont essentielles. Cependant, il \deleted{\unsure{semble}} est possible que la 'hypothèse de quasi-projectivité sur $X$ \deleted{et celle de} \inserted{et $X'$} peut être omise, \deleted{et celle de \unsure{propreté} \unsure{être omise|des variétés|dans les corollaires} dans le Th, et les \unsure{puisse de|premiers cas|partie des}} sauf tout au plus dans cor. 2 (ii) [Existence d'un schéma algébrique donnant un schéma formel].
% archive sheet number: 43
